\section{Primitive elements}
\label{section-primitive-element}

\noindent
Let $E/F$ be a finite extension of fields. An element $\alpha \in E$
is called a {\it primitive element of $E$ over $F$} if $E = F(\alpha)$.

\begin{lemma}[Primitive element]
\label{lemma-primitive-element}
Let $E/F$ be a finite extension of fields. The following are equivalent
\begin{enumerate}
\item there exists a primitive element for $E$ over $F$, and
\item there are finitely many subextensions $E/K/F$.
\end{enumerate}
Moreover, (1) and (2) hold if $E/F$ is separable.
\end{lemma}

\begin{proof}
Let $\alpha \in E$ be a primitive element. Let $P$ be the minimal
polynomial of $\alpha$ over $F$. Next, let $E/K/F$ be a subextension.
Let $Q$ be the minimal polynomial of $\alpha$ over $K$. Observe that
$\deg(Q) = [E : K]$. Writing $Q = x^d + \sum_{i < d} a_i x^i$ we claim that
$K$ is equal to $L = F(a_0, \ldots, a_{d - 1})$. Indeed $\alpha$ has degree
$d$ over $L$ and $L \subset K$. Hence $[E : L] = [E : K]$ and it follows
that $[K : L] = 1$, i.e., $K = L$.
Thus it suffices to show there are at most finitely many possibilities
for the polynomial $Q$. This is clear because we have a factorization
$P = QR$  in $K[x]$ in particular in $E[x]$. Since we have unique
factorization in $E[x]$ there are at most finitely many monic
factors of $P$ in $E[x]$.

\medskip\noindent
If $F$ is a finite field (equivalently $E$ is a finite field), then
$E/F$ has a primitive element by the discussion in
Section \ref{section-finite}.
Next, assume $F$ is infinite and there are at most finitely many proper
subfields $E/K/F$. List them, say $K_1, \ldots, K_N$. Then
each $K_i \subset E$ is a proper sub $F$-vector space. As $F$ is infinite
we can find a vector $\alpha \in E$ with $\alpha \not \in K_i$ for all $i$
(a vector space can never be equal to a finite union of proper subvector spaces;
details omitted). Then $\alpha$ is a primitive element for $E$ over $F$.

\medskip\noindent
Having established the equivalence of (1) and (2) we now turn to
the final statement of the lemma. Choose an algebraic closure
$\overline{F}$ of $F$. Enumerate the elements
$\sigma_1, \ldots, \sigma_n \in \Mor_F(E, \overline{F})$.
Since $E/F$ is separable we have $n = [E : F]$ by
Lemma \ref{lemma-separable-equality}.
Note that if $i \not = j$, then
$$
V_{ij} = \Ker(\sigma_i - \sigma_j : E \longrightarrow \overline{F})
$$
is not equal to $E$. Hence arguing as in the preceding paragraph
we can find $\alpha \in E$ with $\alpha \not \in V_{ij}$ for all
$i \not = j$. It follows that $|\Mor_F(F(\alpha), \overline{F})| \geq n$.
On the other hand $[F(\alpha) : F] \leq [E : F]$. Hence equality
by Lemma \ref{lemma-separable-equality}
and we conclude that $E = F(\alpha)$.
\end{proof}






